%% ============================================================
%%  Nota Metodológica — Impacto Tarifário no Complexo Florestal SC
%%  Observatório FIESC | Junho 2026
%%  Compile com: pdflatex nota_metodologica_impacto_tarifas.tex
%% ============================================================
\documentclass[12pt, a4paper]{article}

%% Pacotes
\usepackage[utf8]{inputenc}
\usepackage[T1]{fontenc}
\usepackage[brazil]{babel}
\usepackage{geometry}
\usepackage{amsmath, amssymb}
\usepackage{booktabs}
\usepackage{array}
\usepackage{multirow}
\usepackage{xcolor}
\usepackage{graphicx}
\usepackage{hyperref}
\usepackage{enumitem}
\usepackage{lmodern}
\usepackage{caption}
\usepackage{float}
\usepackage{siunitx}
\usepackage{makecell}

\geometry{left=3cm, right=2.5cm, top=3cm, bottom=2.5cm}

%% Cores institucionais
\definecolor{azulFIESC}{HTML}{1F4E79}
\definecolor{azulMed}{HTML}{2E75B6}
\definecolor{azulCl}{HTML}{DEEAF1}
\definecolor{verde}{HTML}{375623}
\definecolor{vermelho}{HTML}{C00000}
\definecolor{cinzaTexto}{HTML}{595959}

%% Hyperlinks
\hypersetup{
  colorlinks = true,
  linkcolor  = azulFIESC,
  urlcolor   = azulMed,
  citecolor  = azulFIESC,
  pdftitle   = {Nota Metodológica — Impacto Tarifário no Complexo Florestal SC},
  pdfauthor  = {Observatório FIESC},
}

%% Cabeçalho de seções colorido
\usepackage{titlesec}
\titleformat{\section}{\large\bfseries\color{azulFIESC}}{\thesection}{1em}{}[\color{azulMed}\hrule]
\titleformat{\subsection}{\normalsize\bfseries\color{azulMed}}{\thesubsection}{1em}{}
\titleformat{\subsubsection}{\normalsize\itshape\color{cinzaTexto}}{\thesubsubsection}{1em}{}

%% Configuração de tabelas com siunitx
\sisetup{
  group-separator = {.},
  group-minimum-digits = 4,
  output-decimal-marker = {,},
}

%% ============================================================
\begin{document}

%% ---- CAPA ----
\begin{titlepage}
  \centering
  \vspace*{2cm}
  {\color{azulFIESC}\rule{\linewidth}{2pt}}\\[0.4cm]
  {\LARGE\bfseries\color{azulFIESC}
    Nota Metodológica\\[0.3cm]
    Impacto das Tarifas Americanas sobre\\
    o Complexo Florestal de Santa Catarina
  }\\[0.3cm]
  {\color{azulFIESC}\rule{\linewidth}{1pt}}\\[1.5cm]
  {\large\color{cinzaTexto}
    Observatório FIESC — Competitividade Industrial\\[0.3cm]
    Junho de 2026
  }\\[3cm]
  \begin{minipage}{0.9\textwidth}
    \centering\itshape\color{cinzaTexto}
    Documento técnico de suporte ao relatório SEAFLOR~2026.\\
    Descreve as definições, fórmulas e premissas utilizadas na estimativa\\
    do impacto tarifário sobre as exportações florestais catarinenses\\
    no período de janeiro a maio de 2026.
  \end{minipage}\\[2cm]
  {\footnotesize\color{cinzaTexto}
    Fonte dos dados: ComexStat/MDIC $\cdot$
    Classificação setorial: CNAE/NCM FIESC $\cdot$
    Período de referência: Jan--Mai 2025 vs.\ Jan--Mai 2026
  }
\end{titlepage}

\tableofcontents
\newpage

%% ============================================================
\section{Contexto e Motivação}

Em 2025--2026 o governo dos Estados Unidos reintroduziu uma série de tarifas
sobre importações provenientes do Brasil, afetando em especial os produtos do
Complexo Florestal catarinense (madeira, móveis, papel/celulose e produtos
da base florestal). O presente documento descreve a metodologia empregada
pelo Observatório FIESC para mensurar o impacto dessas tarifas sobre a
receita de exportação de Santa Catarina.

O \textbf{Complexo Florestal} é definido pelos setores com classificação
CNAE correspondente a:

\begin{itemize}[noitemsep]
  \item \textbf{CNAE~2} — Produção Florestal / Base Florestal
  \item \textbf{CNAE~16} — Fabricação de Produtos de Madeira
  \item \textbf{CNAE~17} — Fabricação de Celulose, Papel e Produtos de Papel
  \item \textbf{CNAE~31} — Fabricação de Móveis
\end{itemize}

No plano NCM, essa classificação corresponde aos códigos do SH4 mapeados
no dicionário NCM--FIESC para os capítulos 44, 47, 48 e 94, entre outros.

%% ============================================================
\section{Base de Dados}

\begin{description}[style=nextline]
  \item[Exportações históricas (2024--2025)]
    Arquivo \textit{gold} do ComexStat/MDIC no nível NCM~8 dígitos,
    filtrado para Santa Catarina
    (\texttt{sg\_uf = 'SC'}, \texttt{tp\_carga = 'EXP'}).
  \item[Exportações e importações 2026 (jan--mai)]
    Arquivos bronze mensais \texttt{EXP\_2026.csv} e \texttt{IMP\_2026.csv}
    do ComexStat, enriquecidos pelo dicionário NCM--FIESC
    (tabela \texttt{NCM\_SH4}) e pela tabela de países (\texttt{PAIS.csv}).
  \item[Referência temporal]
    Todos os comparativos utilizam o período \textbf{janeiro a maio}
    (\(m \in \{1,\dots,5\}\)) para garantir a simetria interanual.
\end{description}

Notação adotada ao longo do documento:

\begin{align*}
  X_{d,s}^{t} &= \text{valor FOB exportado (US\$) pelo setor } s
                  \text{ para o destino } d \text{ em Jan--Mai do ano } t\\
  Q_{d,s}^{t} &= \text{volume (kg líquido) exportado pelo setor } s
                  \text{ para o destino } d \text{ em Jan--Mai do ano } t\\
  \text{EUA}  &= \text{destino ``Estados Unidos''}\\
  \text{OUT}  &= \text{conjunto de todos os demais destinos (exceto EUA)}
\end{align*}

%% ============================================================
\section{Metodologia de Decomposição do Impacto Tarifário}

O impacto total das tarifas é decomposto em quatro componentes calculados
para cada setor \(s\) e para o Complexo Florestal como um todo.

%%--------------------------------------------------------------
\subsection{Destruição Bruta}

A \textbf{Destruição Bruta} (\(DB_s\)) mede a queda absoluta nas exportações
do setor \(s\) para os Estados Unidos entre os períodos de referência:

\begin{equation}
  \boxed{
    DB_s = X_{\text{EUA},s}^{2025} - X_{\text{EUA},s}^{2026}
  }
  \label{eq:db}
\end{equation}

Um valor $DB_s > 0$ indica redução nas vendas ao mercado americano.
Um valor negativo (como observado na Base Florestal, CNAE~2) indica que
as exportações para os EUA \emph{cresceram} no período, provavelmente
devido a realocações ou crescimento de nicho não relacionado às tarifas.

%%--------------------------------------------------------------
\subsection{Desvio de Comércio}

O \textbf{Desvio de Comércio} (\(DC_s\)) captura o ganho líquido nas
exportações do setor \(s\) para todos os mercados fora dos EUA,
ou seja, a parcela da perda americana compensada pela diversificação
geográfica. Utiliza-se $\max(\cdot, 0)$ porque uma queda nos demais
mercados não pode compensar: seria perda adicional, não desvio:

\begin{equation}
  \boxed{
    DC_s = \max\!\left(X_{\text{OUT},s}^{2026} - X_{\text{OUT},s}^{2025},\; 0\right)
  }
  \label{eq:dc}
\end{equation}

\noindent\textbf{Atenção:} o desvio é medido em \emph{valor} (US\$ FOB),
não em volume. Como o valor unitário obtido em mercados alternativos é,
em geral, inferior ao praticado nos EUA, o ganho em valor tende a
subestimar o ganho em volume — essa diferença é capturada pelo
Efeito-Preço (seção~\ref{sec:efeito_preco}).

%%--------------------------------------------------------------
\subsection{Destruição Líquida}

A \textbf{Destruição Líquida} (\(DL_s\)) é o impacto líquido após descontar
a compensação obtida nos mercados alternativos:

\begin{equation}
  \boxed{
    DL_s = DB_s - DC_s
  }
  \label{eq:dl}
\end{equation}

Quando \(DL_s < 0\) (caso do Papel e Celulose e da Base Florestal),
o setor \emph{mais do que compensou} a perda com os EUA por meio da
expansão para outros destinos — ainda que possa haver perda de receita
pelo diferencial de preço (ver seção~\ref{sec:efeito_preco}).

%%--------------------------------------------------------------
\subsection{Prêmio de Mercado EUA}
\label{sec:premio}

O mercado americano paga, historicamente, um preço unitário (US\$/kg)
superior ao praticado nos demais destinos. O \textbf{Prêmio de Mercado
EUA} (\(\pi_s\)) é calculado com base nos valores unitários médios de 2025:

\begin{equation}
  VU_{d,s}^{t} = \frac{X_{d,s}^{t}}{Q_{d,s}^{t}}
  \label{eq:vu}
\end{equation}

\begin{equation}
  \boxed{
    \pi_s = VU_{\text{EUA},s}^{2025} - VU_{\text{OUT},s}^{2025}
  }
  \label{eq:premio}
\end{equation}

O prêmio relativo (\%) é dado por:

\begin{equation}
  \pi_s^{\%} = \frac{\pi_s}{VU_{\text{OUT},s}^{2025}} \times 100
  \label{eq:premio_pct}
\end{equation}

A Tabela~\ref{tab:vu} apresenta os valores unitários e prêmios por CNAE
para o ano de 2025 (referência):

\begin{table}[H]
  \centering
  \caption{Valor unitário médio e prêmio de mercado EUA — 2025 (ano completo)}
  \label{tab:vu}
  \renewcommand{\arraystretch}{1.3}
  \begin{tabular}{lS[table-format=1.4]S[table-format=1.4]S[table-format=1.4]S[table-format=4.1]}
    \toprule
    \textbf{Setor / CNAE}
      & \textbf{\makecell{VU EUA\\(US\$/kg)}}
      & \textbf{\makecell{VU Outros\\(US\$/kg)}}
      & \textbf{\makecell{Prêmio\\(US\$/kg)}}
      & \textbf{\makecell{Prêmio\\(\%)}} \\
    \midrule
    Complexo Florestal         & 1.0816 & 0.4654 & 0.6162 & 132.4 \\
    Madeira (CNAE 16)          & 0.9419 & 0.4343 & 0.5075 & 116.8 \\
    Móveis (CNAE 31)           & 2.9810 & 2.1337 & 0.8473 &  39.7 \\
    Papel e Celulose (CNAE 17) & 1.6396 & 0.7565 & 0.8831 & 116.7 \\
    Base Florestal (CNAE 2)    & 4.9854 & 0.0880 & 4.8974 & 5565.9 \\
    \bottomrule
  \end{tabular}
\end{table}

%%--------------------------------------------------------------
\subsection{Efeito-Preço}
\label{sec:efeito_preco}

Quando o setor realoca volume para outros mercados, ele o faz a um preço
unitário inferior ao que obteria nos EUA. O \textbf{Efeito-Preço} (\(EP_s\))
quantifica essa perda de receita por unidade desviada:

\begin{equation}
  \Delta Q_{\text{OUT},s} = Q_{\text{OUT},s}^{2026} - Q_{\text{OUT},s}^{2025}
  \label{eq:delta_kg}
\end{equation}

\begin{equation}
  \boxed{
    EP_s = \max\!\left(\Delta Q_{\text{OUT},s},\; 0\right)
            \times \max\!\left(\pi_s,\; 0\right)
  }
  \label{eq:ep}
\end{equation}

A dupla condição $\max(\cdot, 0)$ garante que:
\begin{enumerate}[noitemsep]
  \item somente o volume \emph{adicional} efetivamente desviado é considerado
        (desvio de volume, não de valor);
  \item somente prêmios positivos geram efeito — setores sem prêmio EUA
        não sofrem perda de preço ao diversificar.
\end{enumerate}

%%--------------------------------------------------------------
\subsection{Perda Total — Metodologia A}
\label{sec:perda_total}

A \textbf{Perda Total} pelo Método~A (\(PT_s^A\)) agrega a destruição
líquida de valor e a perda adicional de receita por diferencial de preço:

\begin{equation}
  \boxed{
    PT_s^A = DL_s + EP_s
  }
  \label{eq:pt}
\end{equation}

Esta é a estimativa base do impacto tarifário. A relação entre os
componentes pode ser resumida como:

\begin{align}
  PT_s^A &= \underbrace{(DB_s - DC_s)}_{\text{destruição líquida}}
            + \underbrace{\Delta Q_{\text{OUT},s}^{+} \cdot \pi_s}_{\text{efeito-preço}}
  \label{eq:pt_expandida}
\end{align}

onde $\Delta Q_{\text{OUT},s}^{+} = \max(\Delta Q_{\text{OUT},s}, 0)$.

%% ============================================================
\section{Análise Contrafactual}
\label{sec:contrafactual}

A análise contrafactual complementa a decomposição ao estimar \emph{quanto
o setor deveria ter exportado para os EUA} caso as tarifas não tivessem
sido impostas, com base em diferentes janelas de referência histórica.

%%--------------------------------------------------------------
\subsection{Formulação}

Para cada cenário \(k\), define-se um \textbf{share EUA de referência}
(\(sh_{\text{EUA},s}^{k}\)), calculado a partir do histórico do setor:

\begin{equation}
  sh_{\text{EUA},s}^{k} = \frac{X_{\text{EUA},s}^{k,\text{ref}}}
                                {X_{\text{TOT},s}^{k,\text{ref}}}
  \label{eq:share_ref}
\end{equation}

O valor \textbf{contrafactual} das exportações para os EUA em 2026 é obtido
aplicando esse share ao total efetivamente exportado:

\begin{equation}
  \hat{X}_{\text{EUA},s}^{k} = X_{\text{TOT},s}^{2026} \times sh_{\text{EUA},s}^{k}
  \label{eq:ctf}
\end{equation}

A \textbf{Receita Forgone} (\(RF_s^k\)) — receita que ``deixou de entrar'' —
é a diferença entre o contrafactual e o valor observado:

\begin{equation}
  \boxed{
    RF_s^k = \hat{X}_{\text{EUA},s}^{k} - X_{\text{EUA},s}^{2026}
  }
  \label{eq:rf}
\end{equation}

%%--------------------------------------------------------------
\subsection{Cenários}

Três horizontes de referência são utilizados, cobrindo desde o ciclo
recente (pós-2025) até o período pré-tarifas Trump:

\vspace{0.5em}
\begin{description}[style=nextline, leftmargin=2cm]
  \item[Cenário A — Conservador]
    $k = \text{Jan--Mai}\ 2025$: baseline mais recente disponível,
    que já reflete alguma antecipação às tarifas. Gera o menor impacto estimado.

  \item[Cenário B — Base]
    $k = 2024\ \text{(ano completo)}$: ano anterior ao choque tarifário,
    consolidado. É o cenário de referência principal.

  \item[Cenário C — Estrutural]
    $k = 2023\ \text{(ano completo)}$: posição pré-ciclo Trump.
    Serve como limite superior e capta a dependência estrutural do setor.
\end{description}

%% ============================================================
\section{Indicadores de Concentração de Mercado}
\label{sec:conc}

O grau de diversificação geográfica das exportações é mensurado por dois
indicadores complementares, calculados sobre a distribuição de valor FOB
por país de destino dentro de cada setor \(s\) e período \(t\).

%%--------------------------------------------------------------
\subsection{Índice Herfindahl-Hirschman (HHI)}

O HHI mede a concentração somando o quadrado das participações percentuais
de cada destino \(i\) no total exportado pelo setor:

\begin{equation}
  \boxed{
    HHI_{s,t} = \sum_{i=1}^{N_{s,t}}
      \left(\frac{X_{i,s}^{t}}{X_{\text{TOT},s}^{t}} \times 100\right)^{2}
  }
  \label{eq:hhi}
\end{equation}

onde $N_{s,t}$ é o número de países destino do setor $s$ no período $t$.

\noindent Classificação adotada (critério DOJ/FTC):

\begin{center}
\renewcommand{\arraystretch}{1.2}
\begin{tabular}{ll}
\toprule
\textbf{HHI} & \textbf{Grau de concentração} \\
\midrule
$< 1.000$    & Mercado diversificado \\
$1.000$--$1.499$ & Baixa concentração \\
$1.500$--$2.499$ & Concentração moderada \\
$\geq 2.500$ & Alta concentração \\
\bottomrule
\end{tabular}
\end{center}

%%--------------------------------------------------------------
\subsection{Índice de Gini}

O Gini de destinos mede a desigualdade na distribuição das exportações
entre países. Sobre o vetor ordenado de forma crescente
$\mathbf{x} = (x_1, x_2, \dots, x_N)$ com $x_1 \leq x_2 \leq \cdots \leq x_N$:

\begin{equation}
  \boxed{
    G_{s,t} = \frac{2 \displaystyle\sum_{i=1}^{N} i\, x_i}
                   {N \displaystyle\sum_{i=1}^{N} x_i}
               - \frac{N+1}{N}
  }
  \label{eq:gini}
\end{equation}

$G = 0$ indica distribuição perfeitamente igualitária (todas as exportações
divididas igualmente entre os destinos); $G = 1$ indica concentração
máxima em um único destino.

%%--------------------------------------------------------------
\subsection{Contribuição individual ao HHI}

Para identificar quais destinos mais pressionam a concentração,
calcula-se a \textbf{contribuição individual de cada país} ao HHI:

\begin{equation}
  HHI_{\text{contrib},i,s}^{t}
    = \left(\frac{X_{i,s}^{t}}{X_{\text{TOT},s}^{t}} \times 100\right)^{2}
  \label{eq:hhi_contrib}
\end{equation}

de forma que $\displaystyle HHI_{s,t} = \sum_{i=1}^{N} HHI_{\text{contrib},i,s}^{t}$.

%% ============================================================
\section{Resultados por CNAE — Jan-Mai 2025 vs.\ 2026}
\label{sec:resultados}

%%--------------------------------------------------------------
\subsection{Decomposição do Impacto Tarifário}

A Tabela~\ref{tab:decomp} apresenta, para cada CNAE do Complexo Florestal,
os cinco componentes da metodologia de decomposição.
Valores em US\$ milhões (FOB). Período: janeiro a maio.

\begin{table}[H]
  \centering
  \caption{Decomposição do impacto tarifário por CNAE — Jan-Mai 2025 vs.\ 2026 (US\$ milhões)}
  \label{tab:decomp}
  \renewcommand{\arraystretch}{1.35}
  \small
  \begin{tabular}{
      l
      S[table-format=3.1]
      S[table-format=2.1]
      S[table-format=3.1]
      S[table-format=2.1]
      S[table-format=3.1]
    }
    \toprule
    \textbf{Setor}
      & \textbf{\makecell{Destruição\\Bruta ($DB_s$)}}
      & \textbf{\makecell{Desvio\\($DC_s$)}}
      & \textbf{\makecell{Destruição\\Líquida ($DL_s$)}}
      & \textbf{\makecell{Efeito-\\Preço ($EP_s$)}}
      & \textbf{\makecell{Perda\\Total ($PT_s^A$)}} \\
    \midrule
    Madeira (CNAE 16)          & 118.1 & 32.5 &  85.6 & 30.1 & 115.8 \\
    Móveis (CNAE 31)           &  24.9 &  0.0 &  24.9 &  0.0 &  24.9 \\
    Papel e Celulose (CNAE 17) &   5.2 & 10.8 &  -5.6 & 17.7 &  12.2 \\
    Base Florestal (CNAE 2)    &  -0.6 &  0.0 &  -0.6 &  0.0 &  -0.6 \\
    \midrule
    \textbf{Complexo Florestal}& \textbf{147.7} & \textbf{29.2}
      & \textbf{118.6} & \textbf{0.0} & \textbf{118.6} \\
    \bottomrule
  \end{tabular}
  \vspace{0.4em}
  \begin{minipage}{\textwidth}
    \footnotesize
    \textit{Notas:} (1) O Efeito-Preço do Complexo Florestal consolidado é zero porque
    o volume total exportado para outros mercados \emph{recuou} no período
    ($\Delta Q_{\text{OUT}} \leq 0$), mesmo com o valor em USD tendo crescido
    — evidência de aumento de preço médio sem expansão de volume.
    (2) Papel e Celulose ($DL < 0$): o setor expandiu suas vendas a outros
    destinos a ponto de superar a queda para os EUA em valor; ainda assim,
    o Efeito-Preço (US\$~17,7~mi) indica perda de receita unitária.
    (3) Base Florestal ($DB < 0$): as exportações para os EUA \emph{cresceram}
    no período, provavelmente por contratos de longo prazo não atingidos pelas tarifas.
  \end{minipage}
\end{table}

%%--------------------------------------------------------------
\subsection{Exposição ao Mercado Americano}

\begin{table}[H]
  \centering
  \caption{Participação dos EUA nas exportações por CNAE — Jan-Mai 2025 vs.\ 2026}
  \label{tab:share}
  \renewcommand{\arraystretch}{1.35}
  \small
  \begin{tabular}{lS[table-format=3.0]S[table-format=3.0]S[table-format=2.2]S[table-format=2.2]S[table-format=-2.2]}
    \toprule
    \textbf{Setor}
      & \textbf{\makecell{EXP EUA 2025\\(US\$ mi)}}
      & \textbf{\makecell{EXP EUA 2026\\(US\$ mi)}}
      & \textbf{\makecell{Share EUA\\2025 (\%)}}
      & \textbf{\makecell{Share EUA\\2026 (\%)}}
      & \textbf{\makecell{Variação\\(p.p.)}} \\
    \midrule
    Madeira (CNAE 16)          & 262 & 144 & 48.58 & 31.72 & -16.86 \\
    Móveis (CNAE 31)           &  48 &  24 & 43.35 & 29.96 & -13.39 \\
    Papel e Celulose (CNAE 17) &  11 &   6 &  7.72 &  4.07 &  -3.65 \\
    Base Florestal (CNAE 2)    &   0 &   1 &  0.47 &  3.94 &  +3.47 \\
    \midrule
    \textbf{Complexo Florestal}& \textbf{322} & \textbf{174}
      & \textbf{39.18} & \textbf{24.79} & \textbf{-14.39} \\
    \bottomrule
  \end{tabular}
\end{table}

%%--------------------------------------------------------------
\subsection{Análise Contrafactual — Receita Foregone}

A Tabela~\ref{tab:ctf} compara os três cenários contrafactuais para o
Complexo Florestal e cada CNAE individualmente.

\begin{table}[H]
  \centering
  \caption{Receita Foregone por cenário contrafactual e CNAE — Jan-Mai 2026 (US\$ milhões)}
  \label{tab:ctf}
  \renewcommand{\arraystretch}{1.35}
  \small
  \begin{tabular}{
      l
      S[table-format=2.2]
      S[table-format=3.1]
      S[table-format=2.2]
      S[table-format=3.1]
      S[table-format=2.2]
      S[table-format=3.1]
    }
    \toprule
    \multirow{2}{*}{\textbf{Setor}}
      & \multicolumn{2}{c}{\textbf{Cenário A} (\textit{conservador})}
      & \multicolumn{2}{c}{\textbf{Cenário B} (\textit{base})}
      & \multicolumn{2}{c}{\textbf{Cenário C} (\textit{estrutural})} \\
    \cmidrule(lr){2-3}\cmidrule(lr){4-5}\cmidrule(lr){6-7}
      & \textbf{\makecell{Baseline\\EUA (\%)}}
      & \textbf{\makecell{RF\\(US\$ mi)}}
      & \textbf{\makecell{Baseline\\EUA (\%)}}
      & \textbf{\makecell{RF\\(US\$ mi)}}
      & \textbf{\makecell{Baseline\\EUA (\%)}}
      & \textbf{\makecell{RF\\(US\$ mi)}} \\
    \midrule
    Madeira (CNAE 16)          & 48.58 &  76.5 & 50.11 &  83.5 & 50.65 &  85.9 \\
    Móveis (CNAE 31)           & 43.35 &  10.5 & 45.68 &  12.3 & 47.64 &  13.9 \\
    Papel e Celulose (CNAE 17) &  7.72 &   5.6 &  7.61 &   5.5 & 10.35 &   9.7 \\
    Base Florestal (CNAE 2)    &  0.47 &  -0.6 &  3.63 &  -0.1 &  3.81 &   0.0 \\
    \midrule
    \textbf{Complexo Florestal}& \textbf{39.18} & \textbf{101.3}
      & \textbf{41.04} & \textbf{114.4}
      & \textbf{41.06} & \textbf{114.5} \\
    \bottomrule
  \end{tabular}
  \vspace{0.4em}
  \begin{minipage}{\textwidth}
    \footnotesize
    \textit{Notas:} Cenário~A: baseline = share EUA observado em Jan--Mai~2025.
    Cenário~B: baseline = share EUA do ano completo de 2024.
    Cenário~C: baseline = share EUA do ano completo de 2023 (pré-ciclo Trump).
    RF = Receita Foregone (\textit{receita que deixou de ser realizada}).
    Valores negativos na Base Florestal (CNAE~2) indicam que as exportações
    para os EUA superaram o contrafactual — o setor ganhou participação americana.
  \end{minipage}
\end{table}

%% ============================================================
\section{Síntese dos Resultados}

\begin{description}[style=nextline]
  \item[Complexo Florestal (consolidado)]
    A queda de 45,8\% nas exportações para os EUA gerou uma
    \textbf{destruição bruta de US\$~147,7~milhões}.
    O desvio de comércio para outros mercados (US\$~29,2~mi) resultou em
    \textbf{destruição líquida de US\$~118,6~milhões}.
    Os cenários B e C convergem em torno de \textbf{US\$~114--115~milhões}
    de receita foregone (14--16\% do total exportado), confirmando a
    robustez do intervalo.

  \item[Madeira (CNAE 16)]
    Concentra \textbf{$\approx$~98\% da destruição bruta total} do complexo
    (US\$~118,1~mi). A dependência dos EUA caiu de 48,6\% para 31,7\%
    (-17~p.p.). O Efeito-Preço (US\$~30,1~mi) evidencia que a madeira
    reorientada a outros destinos perdeu em média US\$~0,50/kg de receita unitária.
    Perda total estimada: US\$~115,8~milhões.

  \item[Móveis (CNAE 31)]
    Perda de US\$~24,9~mi sem nenhum desvio de comércio —
    os mercados alternativos não absorveram volume adicional.
    O alto prêmio EUA (US\$~0,85/kg; 40\% superior) reforça a
    vulnerabilidade do setor na ausência de diversificação.

  \item[Papel e Celulose (CNAE 17)]
    Único setor com \emph{destruição líquida negativa} (−US\$~5,6~mi):
    as exportações para outros mercados cresceram mais que a queda nos EUA.
    Contudo, o Efeito-Preço (US\$~17,7~mi) gerou uma perda total de
    US\$~12,2~mi — evidência de \emph{dumping implícito} para acessar novos destinos.

  \item[Base Florestal (CNAE 2)]
    Impacto desprezível. As exportações para os EUA cresceram no período
    (possivelmente por contratos de longo prazo), resultando em destruição
    bruta negativa (−US\$~0,6~mi).
\end{description}

%% ============================================================
\section*{Referências}
\addcontentsline{toc}{section}{Referências}

\begin{enumerate}[label={[\arabic*]}]
  \item BRASIL. Ministério do Desenvolvimento, Indústria, Comércio e Serviços
        (MDIC). \textit{ComexStat — Sistema de Análise das Informações de
        Comércio Exterior}. Brasília, 2026. Disponível em: \url{https://comexstat.mdic.gov.br}.

  \item FIESC. \textit{Dicionário NCM--CNAE — Mapeamento do Complexo Florestal
        de Santa Catarina}. Versão atualizada em 07/04/2026.
        Florianópolis: Observatório FIESC, 2026.

  \item U.S.\ Department of Justice; Federal Trade Commission.
        \textit{Horizontal Merger Guidelines}, 2010.
        Critérios de classificação HHI adotados na seção~\ref{sec:conc}.

  \item VINER, Jacob. \textit{The Customs Union Issue}. New York: Carnegie
        Endowment for International Peace, 1950.
        Referência conceitual para \textit{trade creation} e \textit{trade diversion}.
\end{enumerate}

%% ============================================================
\appendix

\section{Glossário de Variáveis}
\label{app:glossario}

\begin{center}
\renewcommand{\arraystretch}{1.3}
\begin{tabular}{lp{10cm}}
\toprule
\textbf{Símbolo} & \textbf{Descrição} \\
\midrule
$X_{d,s}^{t}$ & Valor FOB (US\$) das exportações do setor $s$ para o destino $d$ em Jan--Mai do ano $t$ \\
$Q_{d,s}^{t}$ & Volume (kg líquido) exportado pelo setor $s$ para $d$ em Jan--Mai do ano $t$ \\
$VU_{d,s}^{t}$ & Valor unitário médio: $X_{d,s}^{t} / Q_{d,s}^{t}$ (US\$/kg) \\
$DB_s$        & Destruição Bruta: queda nas exportações para os EUA \\
$DC_s$        & Desvio de Comércio: ganho nas exportações para outros mercados \\
$DL_s$        & Destruição Líquida: $DB_s - DC_s$ \\
$\pi_s$       & Prêmio de mercado EUA: $VU_{\text{EUA},s} - VU_{\text{OUT},s}$ (US\$/kg) \\
$EP_s$        & Efeito-Preço: perda de receita por diferencial de preço no volume desviado \\
$PT_s^A$      & Perda Total pelo Método A: $DL_s + EP_s$ \\
$sh_{\text{EUA},s}^{k}$ & Share EUA de referência no cenário contrafactual $k$ \\
$\hat{X}_{\text{EUA},s}^{k}$ & Exportações contrafactuais para os EUA no cenário $k$ \\
$RF_s^k$      & Receita Foregone no cenário $k$: diferença entre contrafactual e observado \\
$HHI_{s,t}$   & Índice Herfindahl-Hirschman de concentração de destinos \\
$G_{s,t}$     & Índice de Gini de concentração de destinos \\
\bottomrule
\end{tabular}
\end{center}

\end{document}
